\section{Equivariant tilting objects and factorizations}

In this section, we prove that equivariant tilting modules induce equivalences of derived factorization categories.

\subsection{Tilting objects and factorizations}\label{general morita}
In this subsection, we prove that a tilting object in a Gorthendieck category induces an equivalence of derived factorization categories. To prove this, we need the following lemmas:


Let $\scrA$ be a Grothendieck category and $(\scrE,\Phi,w)$ a triple as in \eqref{triple}.
\begin{lem}\label{qis}
Assume that we have $\Ext^i_{\scrA}(A,B)=0$ for arbitrary objects $A,B\in \scrE$ and all $i>0$. Then for objects $E,F\in \Fact(\scrE,\Phi,w)$ and an injective resolution $\varepsilon\colon F\to J^{\bullet}$ in $\Ch(\Fact(\scrA,\Phi,w))$, the map
\begin{gather*}
\varepsilon_*\colon\ \Hom_{\dgFact(\scrA,\Phi,w)}(E,F)^p\to\Hom_{\dgFact(\scrA,\Phi,w)}(E,J^{\bullet})^p
\end{gather*}
of cochain complexes of abelian groups is a quasi-isomorphism for every $p\in \bZ$. Here injective resolution means that $\varepsilon$ is a quasi-isomorphism given by an injection $\varepsilon_0\colon F\hookto J^0$, and all components of the factorizations $J^n$ are injective objects in $\scrA$.

\begin{proof}
Since $\Hom_{\dgFact(\scrA,\Phi,w)}(E,F)^p\cong \Hom_{\dgFact(\scrA,\Phi,w)}(E[p],F)^0$, we may assume that $p=0$. If~$\varepsilon\colon F\to J^{\bullet}$ is an injective resolution, then we have the induced injective resolutions $\varepsilon_i\colon F_i\to J^{\bullet}_i$ for $i=0,1$.
Then the map $\varepsilon_*\colon \Hom_{\dgFact(\scrA,\Phi,w)}(E,F)^0\to\Hom_{\dgFact(\scrA,\Phi,w)}(E,J^{\bullet})^0$ of cochain complexes is the direct sum of two maps
\begin{gather*}
(\varepsilon_*)_i\colon\ \Hom_{\scrA}(E_i,F_i)\to \Hom_{\scrA}(E_i,J^{\bullet}_i)\qquad (i=0,1)
\end{gather*}
of cochain complexes. The $i$-th cohomology of the cochain complex $\Hom_{\scrA}(E_i,J^{\bullet}_i)$ is nothing but $\Ext^i_{\scrA}(E_i,F_i)$. Hence, by the assumption, the map $(\varepsilon_*)_i$ is a quasi-isomorphism, and so is~$\varepsilon_*$.
\end{proof}
\end{lem}


\begin{lem}\label{add ff}
Assume that we have $\Ext^i_{\scrA}(A,B)=0$ for arbitrary objects $A,B\in \scrE$ and all $i>0$. Then the natural functor $\K(\scrE,\Phi,w)\to \Dco(\scrA,\Phi,w)$ is fully faithful.
\begin{proof}
Let $E,F\in \Fact(\scrE,\Phi,w)$ be objects. We need to show that the map
\begin{gather*}
Q\colon\Hom_{{\rm K}(\scrA,\Phi,w)}(E,F)\to \Hom_{\Dco(\scrA,\Phi,w)}(E,F)
\end{gather*}
 defined by the Verdier quotient $Q\colon{\rm K}(\scrA,\Phi,w)\to\Dco(\scrA,\Phi,w)$ is an isomorphism. By~\cite[Proposition~2.19]{bdfik} we can take an injective resolution $\iota\colon F\to (I^{\bullet},d_{I}^{\bullet})=\big(\cdots\to0\to I^0\xrightarrow{d_{I}^0} I^1\xrightarrow{d_{I}^1} \cdots\big)$ of $F$ in $\Ch(\Fact(\scrA,\Phi,w))$ in the sense that all $I^i$ has injective components. Then, we have the induced map
 \begin{gather*}
 \Tot(\iota)\colon\ F\to\Tot(I^{\bullet})
 \end{gather*}
 in $\Fact(\scrA,\Phi,w)$ which becomes an isomorphism in $\Dco(\scrA,\Phi,w)$. Consider the following commutative diagram:
\[
\begin{tikzcd}
\Hom_{{\rm K}(\scrA,\Phi,w)}(E,F)\arrow[rr, "Q"]\arrow[d, "\Tot(\iota)_*"']&&\Hom_{\Dco(\scrA,\Phi,w)}(E,F)\arrow[d, "\Tot(\iota)_*"]\\
\Hom_{{\rm K}(\scrA,\Phi,w)}(E,\Tot(I^{\bullet}))\arrow[rr, "Q"]&&\Hom_{\Dco(\scrA,\Phi,w)}(E,\Tot(I^{\bullet})).
\end{tikzcd}
\]
In the above diagram, the vertical arrow on the right hand side is an isomorphism, and the horizontal arrow on the bottom is also an isomorphism by~\cite[Lemma~2.24]{bdfik} and~\cite[Proposition~B.2(I)]{ls} (see also~\cite[Remark 2.14]{ls}). Hence it suffices to show that the vertical arrow on the left hand side is an isomorphism, and for this we prove that the map
\begin{equation}\label{hom cpx}
\Tot(\iota)_*\colon\ \Hom_{\dgFact(\scrA,\Phi,w)}(E,F) \to\Hom_{\dgFact(\scrA,\Phi,w)}(E,\Tot(I^{\bullet}))
\end{equation}
of cochain complexes is a quasi-isomorphism. The complex $\Hom_{\dgFact(\scrA,\Phi,w)}(E,F)$ can be seen as the double complex $X^{\bullet,\bullet}$ defined by
\begin{equation*}
X^{p,q}\defeq\begin{cases}
\Hom_{\dgFact(\scrA,\Phi,w)}(E,F)^p, & q=0, \\
0,& q\neq 0.
\end{cases}
\end{equation*}
On the other hand, we can consider the double complex $Y^{\bullet,\bullet}$ defined by
\begin{equation*}
Y^{p,q}\defeq\Hom_{\dgFact(\scrA,\Phi,w)}(E,I^q)^p,
\end{equation*}
 where the differentials $d_Y^{p,\bullet}\colon Y^{p,\bullet}\to Y^{p,\bullet+1}$ and $d_Y^{\bullet,q}\colon Y^{\bullet,q}\to Y^{\bullet+1,q}$ are given by $d_Y^{p,\bullet}\defeq (d_I^{\bullet})_*$ and $d_Y^{\bullet,q}\defeq d_{(E,I^q)}^{\bullet}$, where $d_{(E,I^q)}^{\bullet}$ is the differential of the complex $\Hom_{\dgFact(\scrA,\Phi,w)}(E,I^{q})^{\bullet}$. Then the injective map $\iota\colon F\to I^0$ defining the resolution $\iota\colon F\to I^{\bullet}$ gives the morphism
 \begin{equation}\label{hom cpx 2}
 \iota_*\colon\ X^{\bullet,\bullet}\to Y^{\bullet,\bullet}
 \end{equation}
 of double complexes. By definition, we have isomorphisms $\Tot(X^{\bullet,\bullet})=\Hom_{\dgFact(\scrA,\Phi,w)}(E,F)$ and $\Tot(Y^{\bullet,\bullet})\cong\Hom_{\dgFact(\scrA,\Phi,w)}(E,\Tot(I^{\bullet}))$, and the map $\Tot(\iota)_*$ in (\ref{hom cpx}) is the totalization $\Tot(\iota_*)$ of the map $\iota_*$ in (\ref{hom cpx 2}).
By Lemma~\ref{qis} the map $\iota_*\colon X^{p,\bullet}\to Y^{p,\bullet}$ of cochain complexes is a quasi-isomorphism for any $p\in \bZ$, and therefore the map in (\ref{hom cpx}) is also a quasi-isomorphism by~\cite[Theorem~1.9.3]{ks} (see also~\cite[Lemma~2.46]{ls}). This completes the proof.
\end{proof}
\end{lem}

\begin{lem}\label{canonical truncation}
Let $F^{\bullet}\in \D^+(\scrA)$ be an object, and let $m\in \bZ$ be an integer such that $F^k=0$ for any $k<m$. Then there are a family $\{F^{\bullet}_i\}_{i\in \bZ}$ of objects in $\Db(\scrA)$ with $F^{k}_i=0$ for any $k<m$ and $i$ and an exact triangle
\begin{gather*}
\bigoplus_{i\in \bZ}F^{\bullet}_i\to \bigoplus_{i\in \bZ}F^{\bullet}_i\to F^{\bullet}\to \bigoplus_{i\in \bZ}F^{\bullet}_i[1].
\end{gather*}
\begin{proof}
Denote by $\tau_{\leq i}(F^{\bullet})$ the canonical truncation of $F^{\bullet}$ defined by $\tau_{\leq i}(F^{\bullet})^j=F^j$ if $j<i$, $\tau_{\leq i}(F^{\bullet})^i=\Ker(d^i\colon F^i\to F^{i+1})$ and $\tau_{\leq i}(F^{\bullet})^j=0$ if $j>i$. If we set $F_i^{\bullet}\defeq \tau_{\leq i}(F^{\bullet})$, $F_i^{\bullet}$ lies in~$\Db(\scrA)$ and $F^{k}_i=0$ for any $k<m$ and $i$. The object $F^{\bullet}$ is the cokernel of an injective morphism
\begin{gather*}
\bigoplus_{i\in \bZ} F^{\bullet}_i\to \bigoplus_{i\in \bZ} F^{\bullet}_i
\end{gather*}
in the abelian category $\Ch(\scrA)$ given by
 $\id \oplus (-\iota_i)\colon F^{\bullet}_i\to F^{\bullet}_i\oplus F^{\bullet}_{i+1}$, where $\iota_i\colon F^{\bullet}_i\to F^{\bullet}_{i+1}$ is a~natural injection. Hence we have a triangle
\begin{gather*}
\bigoplus_{i\in \bZ}F^{\bullet}_i\to \bigoplus_{i\in \bZ}F^{\bullet}_i\to F^{\bullet}\to \bigoplus_{i\in \bZ}F^{\bullet}_i[1]
\end{gather*}
in $\D^+(\scrA)$.
\end{proof}
\end{lem}

Let $\scrB$ be a Grothendieck category with enough projectives, and let $F\colon \scrA\to\scrB$ be an addi\-tive functor such that it commutes with small direct sums and has a left adjoint functor $G\colon \scrB\to \scrA$. The adjunction $G\dashv F$ implies that $F$ is left exact and $G$ is right exact. By~\cite[Proposition~14.3.4]{ks2}, $F$ admits a right derived functor
 \begin{gather*}
 \bR F\colon\ \D(\scrA)\to \D(\scrB),
 \end{gather*}
 which commutes with small direct sums, and $ \bR F$ restricts to a functor $\bR^+\!F\colon \D^+(\scrA)\to \D^+(\scrB)$. On the other hand, by~\cite[Theorem~14.4.3]{ks2}, $G$ admits a left derived functor
 \begin{gather*}
 \bL G\colon\ \D(\scrB)\to \D(\scrA),
 \end{gather*}
 and by~\cite[Theorem~14.4.5]{ks2} the adjunction $G\dashv F$ induces an adjunction
 \begin{gather*}
 \bL G\dashv \bR F.
 \end{gather*}

\begin{lem}\label{above ff}
 Assume that $\bR^+\!F\colon \D^+(\scrA)\to \D^+(\scrB)$ restricts to the functor $\bR F\colon \Db(\scrA)\to\Db(\scrB)$ and that there exists a positive integer $d>0$ such that for every object $B\in \scrB$ the projective dimension of $B$ is less than $d$.
 Then, we have the following:
 \begin{itemize}\itemsep=0pt
 \item[$1.$] The functor $\bL G$ restricts to a functor $\bL^{+}G\colon \D^+(\scrB)\to \D^+(\scrA)$.
 \item[$2.$] If $\bR F\colon \Db(\scrA)\to\Db(\scrB)$ is fully faithful, so is $\bR^+\!F\colon \D^+(\scrA)\to \D^+(\scrB)$.
 \end{itemize}
\begin{proof}
(1)
Let $B^{\bullet}\in \D^+(\scrB)$ be an object, and let $m\in\bZ$ be an integer such that $B^k=0$ for any $k<m$. Then by Lemma~\ref{canonical truncation} there are a family $\{B^{\bullet}_i\}_{i\in\bZ}$ of objects in $\Db(\scrB)$ with $B_i^k=0$ for any $k<m$ and $i$ and an exact sequence
\begin{equation}\label{truncation triangle}
\bigoplus_{i\in \bZ}B^{\bullet}_i\to \bigoplus_{i\in \bZ}B^{\bullet}_i\to B^{\bullet}\to \bigoplus_{i\in \bZ}B^{\bullet}_i[1].
\end{equation}
We set $C_i^{\bullet}\defeq\bL G(B^{\bullet}_i)$. Then, by the assumption, $C_i^{k}=0$ for any $k<m-d$ and $i$, and so $\bigoplus_{i\in \bZ}C_i^{\bullet}$ lies in $\D^+(\scrA)$.
Since $\bL G$ admits a right adjoint functor, it commutes with small direct sums.
Applying the functor $\bL G$ to the triangle \eqref{truncation triangle}, we obtain a triangle
\begin{gather*}
\bigoplus_{i\in \bZ}C^{\bullet}_i\to \bigoplus_{i\in \bZ}C^{\bullet}_i\to \bL G(B^{\bullet})\to \bigoplus_{i\in \bZ}C^{\bullet}_i[1],
\end{gather*}
which implies that $\bL G(B^{\bullet})\in\D^+(\scrA)$.


(2) Since $\bR^+\!F$ admits a left adjoint functor $\bL^+G$ by (1), it is enough to show that the adjunction morphism
\begin{gather*}
\varepsilon\colon\ \Upphi\defeq \bL^+G\circ \bR^+\!F\to \id_{\D^+(\scrA)}
\end{gather*}
is an isomorphism of functors. Let $A^{\bullet}\in \D^+(\scrA)$ be an object. Then by Lemma~\ref{canonical truncation} there are a family $\{A^{\bullet}_i\}_{i\in\bZ}$ of objects in $\Db(\scrA)$ and an exact triangle
\begin{gather*}
\bigoplus_{i\in \bZ}A^{\bullet}_i\xrightarrow{f} \bigoplus_{i\in \bZ}A^{\bullet}_i\xrightarrow{g} A^{\bullet}\to \bigoplus_{i\in \bZ}A^{\bullet}_i[1].
\end{gather*}
Consider the following commutative diagram
\[
\begin{tikzcd}
\Upphi\big(\bigoplus_{i\in \bZ}A^{\bullet}_i\big)\arrow[rr, "\Upphi(f)"]\arrow[d, "\varepsilon"']&&\Upphi\big(\bigoplus_{i\in \bZ}A^{\bullet}_i\big)\arrow[rr, "\Upphi(g)"]\arrow[d, "\varepsilon"']&&\Upphi(A^{\bullet})\arrow[d, "\varepsilon"]
\\
\bigoplus_{i\in \bZ}A^{\bullet}_i\arrow[rr, "f"]&&\bigoplus_{i\in \bZ}A^{\bullet}_i\arrow[rr, "g"]&&A^{\bullet},
\end{tikzcd}
\]
where the horizontal sequences are triangles and the vertical arrows are induced by $\varepsilon$. Since $\Upphi$ commutes with direct sums, the vertical arrows on the left and middle are the direct sums of morphisms $\varepsilon_i\defeq\varepsilon(A_i^{\bullet})\colon\Upphi(A_i^{\bullet})\to A_i^{\bullet}$, and each $\varepsilon_i$ is an isomorphism since $\bR F\colon \Db(\scrA)\to\Db(\scrB)$ is fully faithful and $A_i^{\bullet}\in\Db(\scrA)$. Hence the vertical arrow on the right hand side is also an~iso\-morphism.
\end{proof}
\end{lem}


Now we are ready to prove that a tilting object induces an equivalence of derived factorization categories. Let $T\in \scrA$ be a partial tilting object such that the exact subcategory $\Add T\subset \scrA$ is preserved by $\Phi$. Set $\Lambda\defeq\End_{\scrA}(T)$, and
 let $(\Psi,v)$ be a potential on $\Mod\Lambda$. We define
 \begin{gather*}
 \DcoMod(\Lambda,\Psi,v)\defeq\Dco(\Mod\Lambda,\Psi,v),
 \\
 \Dmod(\Lambda,\Psi,v)\defeq\D(\fmod\Lambda,\Psi,v).
 \end{gather*}
Consider the functor
 \begin{gather*}
 \F\defeq\Hom_{\scrA}(T,-)\colon\ \scrA\to \Mod\Lambda,
 \end{gather*}
 and assume that there is a left adjoint functor $\G\colon \Mod\Lambda\to \scrA$, and that $\F$ and $\G$ are factored with respect to $(\Phi,w)$ and $(\Psi,v)$. Let
 \begin{gather*}
 \scrC\subset \scrA
 \end{gather*}
 be an abelian subcategory preserved by $\Phi$ such that $\add T\!\subset \scrC$, $\F(\scrC)\!\subset \fmod\Lambda$ and \mbox{$\G(\fmod\Lambda)\subset \scrC$}. Assume that the natural functors $\Db(\scrC)\to\Db(\scrA)$ and $\D(\scrC,\Phi,w)\to \Dco(\scrA,\Phi,w)$ are fully faithful and that $\bR \F\colon\Db(\scrA)\to\Db(\Mod\Lambda)$ restricts to $\bR \F\colon\Db(\scrC)\to\Db(\fmod\Lambda)$. Then $\bR \F\colon\Dco(\scrA,\Phi,w)\to \DcoMod(\Lambda,\Psi,v)$ restricts to the functor
\begin{gather*}
\bR \F\colon\ \D(\scrC,\Phi,w)\to \Dmod(\Lambda,\Psi,v).
\end{gather*}
Moreover we assume that $\Lambda$ is of finite global dimension. Then we have the left derived functor
\begin{gather*}
\bL\G\colon\ \Dmod(\Lambda,\Psi,v)\to \D(\scrC,\Phi,w)
\end{gather*}
that is left adjoint to $\bR\F\colon\D(\scrC,\Phi,w)\to \Dmod(\Lambda,\Psi,v)$.
\begin{thm}\label{main thm}

Suppose that the adjunction morphisms $\Lambda\to \F\G(\Lambda)$ and $\G\F(T)\to T$ are isomorphisms.
\begin{itemize}\itemsep=0pt
\item[$1.$] The functor $\bL \G\colon \Dmod(\Lambda,\Psi,v)\to \D(\scrC,\Phi,w)$ is fully faithful, and it factors through an equivalence
\begin{gather*}
\bL \G\colon \Dmod(\Lambda,\Psi,v)\simto \Kadd(T,\Phi,w).
\end{gather*}

\item[$2.$]Assume that every object in $\scrA$ has finite injective resolution and that the right derived functor $\bR \F\colon\Db(\scrA)\to\Db(\Mod\Lambda)$ is an equivalence. Then the functor
\begin{gather*}
\bR \F\colon\ \Dco(\scrA,\Phi,w)\to \DcoMod(\Lambda,\Psi,v)
\end{gather*}
is an equivalence, and it restricts to the equivalence
\begin{gather*}
\bR \F\colon\ \D(\scrC,\Phi,w)\simto \Dmod(\Lambda,\Psi,v).
\end{gather*}
\end{itemize}
\begin{proof}
(1) Recall that the functor $\bL\G\colon\Dmod(\Lambda,\Psi,v)\to \D(\scrC,\Phi,w)$ is the composition
\[
\begin{tikzcd}
\Dmod(\Lambda,\Psi,v)\arrow[r,"\sim"]&\K(\proj\Lambda,\Psi,v)\arrow[r,"\G"] &\Kadd(T,\Phi,w)\arrow[r,"Q"]&\D(\scrC,\Phi,w),
\end{tikzcd}
\]
where the first equivalence follows from Proposition~\ref{proj resol} and $Q$ is the natural quotient functor. Since we have the equivalence $\G\colon \proj\Lambda\simto \add T$, the middle functor $\G$ is also an equivalence. Thus the statement follows from Lemma~\ref{add ff} and the assumption that the natural functor $\D(\scrC,\Phi,w)\to \Dco(\scrA,\Phi,w)$ is fully faithful.


%%%%%%%%
(2) Since $(\bR \F)^+\colon\D^+(\scrA)\to\D^+(\Mod\Lambda)$ is fully faithful by Lemma~\ref{above ff}, the functor $\bR \F $: $\Dco(\scrA,\Phi,w)\to \DcoMod(\Lambda,\Psi,v)$ is also fully faithful by~\cite[Lemma~4.11]{bdfik}. Consider the following commutative diagram:
\[
\begin{tikzcd}
\K\Add(T,\Phi,w)\arrow[rr, "\F"]\arrow[d, ]&&\K\Proj(\Lambda,\Psi,v)\arrow[d, ]\\
\Dco(\scrA,\Phi,w)\arrow[rr, "\bR \F"]&&\DcoMod(\Lambda,\Psi,v).
\end{tikzcd}
\]
Since the additive functor $\F\colon\Add(T)\to \Proj\Lambda$ is an equivalence, the top horizontal arrow is an equivalence. Moreover, the vertical arrow on the right hand side is essentially surjective by~Proposition~\ref{proj resol}. Hence $\bR \F\colon\Dco(\scrA,\Phi,w)\to \DcoMod(\Lambda,\Psi,v)$ is essentially surjective, and so it is an equivalence.

Since the natural embedding functors $\D(\scrC,\Phi,w)\to \Dco(\scrA,\Phi,w)$ and $\Dmod(\Lambda,\Psi,v)\to \DcoMod(\Lambda,\Psi,v)$ are fully faithful, the functor $\bR \F\colon\!\D(\scrC,\Phi,w)\!\to \Dmod(\Lambda,\Psi,v)$ and its left~adjo\-int $\bL \G\colon \! \Dmod(\Lambda,\Psi,v)\to \D(\scrE,\Phi,w)$ are fully faithful. Hence $\bR \F\colon\!\D(\scrC,\Phi,w)\to
\Dmod(\Lambda,\Psi,v)$ is also an equivalence.
\end{proof}
\end{thm}

\subsection{Derived factorization categories of noncommutative gauged LG models}

In this short subsection, we define (noncommutative) gauged Landau--Ginzburg models, and its derived factorization categories.

\begin{dfn}
We call the data $(\cA,L,w)^G$ a {\it gauged Landau--Ginzburg model} when $G$ is an~alge\-braic group acting on a scheme $X$, $\cA$ is a $G$-equivariant coherent $X$-algebra, $L$ is a $G$-equivariant line bundle on $X$, and $w\in \Gamma(X,L)^G$ is a $G$-invariant global section of $L$. A gauged Landau--Ginzburg model $(\cA,L,w)^G$ is said to be {\it noncommutative}, if the algebra $\cA$ is not commutative. We write $(X,L,w)^G\defeq(\cO_X,L,w)^G$, and for a character $\chi\colon G\to \bG_m$ of $G$ we write $(\cA,\chi,w)^G\defeq(\cA,\cO(\chi),w)^G$.
\end{dfn}

If $(\cA,L,w)^G$ is a gauged LG model, we have the triple
\begin{gather*}
\big(\Qcoh_G\cA,L\otimes_{\cO_X}(-),w\big)
\end{gather*}
as in \eqref{triple}, where $L\otimes_{\cO_X}(-)\colon \Qcoh_G\cA\simto \Qcoh_G\cA$ is the tensor product with $L$, and $w\colon\id\to L\otimes_{\cO_X}(-)$ is the functor morphism defined by the multiplication by $w$. Then we define
\begin{gather*}
\Qcoh_G(\cA,L,w)\defeq\Fact\big(\Qcoh_G\cA,L\otimes_{\cO_X}(-),w\big),
\\[.5ex]
\DcoQcoh_G(\cA,L,w)\defeq\Dco\big(\Qcoh_G\cA,L\otimes_{\cO_X}(-),w\big),
\\[.5ex]
\Dcoh_G(\cA,L,w)\defeq\D\big(\coh_G\cA,L\otimes_{\cO_X}(-),w\big),
\\[.5ex]
\DMF_G(\cA,L,w)\defeq\D\big(\lfr_G\cA,L\otimes_{\cO_X}(-),w\big),
\end{gather*}
{\sloppy
where $\lfr_G\cA$ is the subcategory of $\coh_G\cA$ consisting of locally free $\cA$-modules. We call $\Dcoh_G(\cA,L,w)$ the {\it derived factorization category} of $(\cA,L,w)^G$, and $\DMF_G(\cA,L,w)$ the {\it derived matrix factorization category} of $(\cA,L,w)^G$. If $X=\Spec R$ is an affine scheme and $A\defeq\Gamma(X,\cA)$ is the corresponding $G$-equivariant $R$-algebra, we define
\begin{gather*}
\Mod_G(A,L,w)\defeq\Fact\big(\Mod_GA,L\otimes_{R}(-),w\big),
\\[.5ex]
\DcoMod_G(A,L,w)\defeq\Dco\big(\Mod_G\cA,L\otimes_{R}(-),w\big),
\\[.5ex]
\Dmod_G(A,L,w)\defeq\D\big(\fmod_GA,L\otimes_{R}(-),w\big),
\\[.5ex]
\KMF_G(A,L,w)\defeq\K\big(\proj_GA,L\otimes_{R}(-),w\big).
\end{gather*}

}


\subsection{Equivariant tilting modules and factorizations}
In this subsection, we apply the general result in Section~\ref{general morita} to the following geometric setting.

Let the notation be the same as in Section~\ref{derived section}.
Let $\chi\colon G/H\to \bG_m$ be a character of $G/H$, and set $\widehat{\chi}\defeq\chi\circ p\colon G\to \bG_m$, where $p\colon G\to G/H$ is the natural projection. Write for $\cO_R(\chi)$ (resp.\ $\cO_X\big(\widehat{\chi}\big)$) the associated $G/H$-equivariant invertible sheaf on $\Spec R$ (resp.\ $G$-equivariant invertible sheaf on $X$). Take a $G/H$-invariant global section $w_R\in \Gamma(\Spec R,\cO_R(\chi))^{G/H}$ of~$\cO_R(\chi)$, and set $w\defeq (f)^*_{p}w_R\in \Gamma\big(X,\cO_X\big(\widehat{\chi}\big)\big)^G$, where $(f)^*_{p}\colon \Mod_{G/H}R\to \Qcoh_GX$ is the pull-back by the $p$-equivariant morphism $f$.
Then we have the gauged LG models
\begin{gather*}
\big(\cA,\widehat{\chi},w\big)^G \qquad \text{and}\qquad \big(\Lambda^H,\chi,w_R\big)^{G/H},
\end{gather*}
and we define
\begin{gather*}
\Kadd_G\big(\cT,\widehat{\chi},w\big)\defeq \K\big(\add_{\widehat{\chi}}\cT,\cO\big(\widehat{\chi}\big)\otimes_{\cO_X}(-),w\big),
\end{gather*}
where $\add_{\widehat{\chi}}\cT$ is the subcategory of $\coh_G\cA$ defined by the following additive closure
\begin{gather*}
\add_{\widehat{\chi}}\cT\defeq\add\big\{\cO\big(\widehat{\chi}^n\big)\otimes_{\cO_X} \big(\cT,\uptheta^{\cT}\big)\mid n\in \bZ \big\}.
\end{gather*}


Since the functors $\F\colon \Qcoh_{G}\cA\to \Mod_{G/H}\Lambda^H$ and $\G\colon \Mod_{G/H}\Lambda^H\to \Qcoh_{G}\cA$ are factored with respect to $(\widehat{\chi},w)$ and $(\chi,w_R)$, they induce the functors
\begin{gather*}
\F\colon\ \Qcoh_G\big(\cA,\widehat{\chi},w\big)\to\Mod_{G/H}\big(\Lambda^H,\chi,w_R\big),
\\
\G\colon\ \Mod_{G/H}\big(\Lambda^H,\chi,w_R\big)\to \Qcoh_G\big(\cA,\widehat{\chi},w\big).
\end{gather*}
Since $\Qcoh_G\cA$ has enough injectives, we have the right derived functor
\begin{gather*}
\bR\F\colon\ \DcoQcoh_G\big(\cA,\widehat{\chi},w\big)\to\DcoMod_{G/H}\big(\Lambda^H,\chi,w_R\big),
\end{gather*}
and it restricts to the functor
\begin{gather*}
\bR\F\colon\ \Dcoh_G\big(\cA,\widehat{\chi},w\big)\to\Dmod_{G/H}\big(\Lambda^H,\chi,w_R\big).
\end{gather*}
If $G/H$ is reductive and every $\Lambda^H$-module has a finite projective resolution, then every finitely generated $G/H$-equivariant $\Lambda^H$-module has a finite projective resolution in $\fmod_{G/H}\Lambda^H$ by Lemma~\ref{finite proj resol}. In this case, we have the left derived functor
\begin{gather*}
\bL\G\colon\ \Dmod_{G/H}\big(\Lambda^H,\chi,w_R\big)\to \Dcoh_{G}\big(\cA,\widehat{\chi},w\big).
\end{gather*}

\begin{thm}\label{tilt and fact}
Assume that $G/H$ is reductive and $\Lambda^H$ is of finite global dimension.
\begin{itemize}\itemsep=0pt
\item[$1.$] If $\big(\cT,\uptheta^{\cT}\big)$ is partial $(G,H)$-tilting, then the functor
\begin{gather*}
\bL\G\colon\ \Dmod_{G/H}\big(\Lambda^H,\chi,w_R\big)\to \Dcoh_{G}\big(\cA,\widehat{\chi},w\big)
\end{gather*}
is fully faithful, and it restricts to the equivalence
\begin{gather*}
\bL\G\colon\ \Dmod_{G/H}\big(\Lambda^H,\chi,w_R\big)\simto \Kadd_{G}\big(\cT,\widehat{\chi},w\big).
\end{gather*}
\item[$2.$] If $\big(\cT,\uptheta^{\cT}\big)$ is $(G,H)$-tilting and every $\cM\in \Qcoh\cA$ has a finite injective resolution in $\Qcoh\cA$, then we have the equivalence
\begin{gather*}
\bR\F\colon\ \Dcoh_{G}\big(\cA,\widehat{\chi},w\big)\simto \Dmod_{G/H}\big(\Lambda^H,\chi,w_R\big).
\end{gather*}
\end{itemize}
\begin{proof}
By an identical argument as in the proof of Lemma~\ref{faithful lemma}, we see that the restriction functors
\begin{gather*}
\Res^G_H\colon\ \Dcoh_G\big(\cA,\widehat{\chi},w\big)\to \Dcoh_H(\cA,\chi,w),
\\
\Res\colon\ \Dmod_{G/H}\big(\Lambda^H,\chi,w_R\big)\to \Dmod\big(\Lambda^H,\id,w_R\big)
\end{gather*}
are faithful. Let $\F_H\colon \coh_H\cA\to\fmod\Lambda^H$ and $\G_H\colon\fmod\Lambda^H\to\coh_H\cA$
be the same functors as in the proof of Theorem~\ref{derived equiv}, and recall that the adjunction morphisms
$\Lambda^H\to \F_H\G_H\big(\Lambda^H\big)$ and $\G_H\F_H(\cT_H)\to \cT_H$
are isomorphisms. Then these functors define the functors
\begin{gather*}
\bR\F_H\colon\ \Dcoh_H\big(\cA,\widehat{\chi},w\big)\to \Dmod\big(\Lambda^H,\chi,w_R\big),
\\
\bL\G_H\colon\ \Dmod\big(\Lambda^H,\chi,w_R\big)\to \Dcoh_H\big(\cA,\widehat{\chi},w\big),
\end{gather*}
and if $\big(\cT,\uptheta^{\cT}\big)$ is partial $(G,H)$-tilting (resp.\ $(G,H)$-tilting), $\bL\G_H$ is fully faithful (resp.\ an~equi\-valence) by Theorem~\ref{main thm}.
Hence the results follow from Lemma~\ref{reduction lemma}.
\end{proof}
\end{thm}
